\chapter{Energiekalibration}
Zur Energiekalibration werden an die Aufgenommenen Energiespektren Gaußkurven mit konstatem Untergrund angepasst. Bei Sich überlagernden Linien werden bis zu zwei Gaußkurven gleichzeitig angepasst.

Die Gaußkurven haben dabei die Form

\begin{align}
	f(x) = \frac{A}{\sqrt{2 \pi} \sigma} \exp(\frac{- (x - \mu)^2}{2 \sigma^2}) + c
\end{algin}

\subsection*{Linke Seite}
\todo{cite k alpha linie source}
Die Anpassungsparameter der Photopeaks sind in \tab{links_na_ganz} und \tab{links_ba_ganz} zu finden. Diesen Linien wurden nun die Relevanten Linien der jeweiligen Elemente zugeordnet. Dabei wurden relative Lagen und Intensitäten betrachtet. Die Zuordnung ist in \tab{links_energy_cal.table} zu finden. Diese Energien und Positionen der Linien werden nun gegeneinander aufgetragen und daran eine Gradengleichung der Form $\mu(E) = a \cdot E + b$ angepasst. zu sehen ist dies in \fig{links/energy_cal}. Für diese Anpassung wurde \texttt{scipy curve\_fit} genutzt. Die so gefundenen Anpassungsparameter sind $a = 1$ und $b = 2$ \todo{add values}. Als Maß der Qualität dieser Anpassungen wurde die Anpassungsgüte verwendet, da das reduzierte Chi Quadrat nur bei Distributionen anwendbar ist, was eine Gradengleichung nicht erfüllt. \todo{cite this} Die Anpassungsgüten von nahe 1 legen nahe, dass eine gute Linearität in der Energiedomäne vorliegt und die Linien korrekt zugeordnet wurden.
\jafps{links/na_ganz}{Mit der linken Seite Aufgenommenes \na Spektrum mit angepassten Gaußkurven}{0.75}
\jafps{links/ba_ganz}{Mit der linken Seite Aufgenommenes \ba Spektrum mit angepassten Gaußkurven}{1}
\jafps{links/energy_cal}{}{0.75}
\jaftl{|c|c|c|c|c|}{../data/fits/links/na_ganz.table}{naganz}{links_na_ganz}
\jaftl{|c|c|c|c|c|}{../data/fits/links/ba_ganz.table}{baganz}{links_ba_ganz}
\jaft{|c|c|}{links_energy_cal.table}{\todo{caption this}}

\todo{energieauflösung}

Die genaue Lage der SCA Schwellen wird aus dem Spektrum mit eingestellten Schwellen entnommen, indem an das Spektrum eine Gaußkurve multipliziert mit zwei Errorfunktionen, eine für die steigende, eine für die fallende Kante, angepasst wird. Die funktion, welche angepasst wird ist dann

\begin{align}
	f(x, x_0, x_1, a_0, a_1, A, \mu, \sigma) &= k(x, x_0, a_0) \cdot (g(x, A, \mu, \sigma) \cdot h(x, x_1, a_1) + c)\\
	k(x, x_0, a_0) &= 0.5 \cdot (1 + \erf(a_0 (x - x_0))) \\
	h(x, x_1, a_1) &= 0.5 \cdot (1 - \erf(a_1 (x - x_1))) \\
	g(x, A, \mu, \sigma) &= \frac{A}{\sqrt{2 \pi} \sigma} \exp(\frac{- (x - \mu)^2}{2 \sigma^2})
\end{align}

Hierbei dient die 

\todo{energie der SCA schwellen}
\todo{energieintervalle der sca fenster}


\subsection*{Rechte Seite}
Die wurde für die rechte Seite wiederholt. Die Spektren sind in \fig{rechts/na_ganz} und \fig{rechts/ba_ganz}, die Anpassungsparameter der Linien in \tab{rechts_na_ganz} und \tab{rechts_ba_ganz} und die Zuordnungen der Linien zu den Energien sind in \tab{rechts_energy_cal.table}. Die Gradengleichung wird wie links angepasst und ergibt die Anpassungsparameter $a = 1$, $b = 2$. \todo{werte}
\jafps{rechts/na_ganz}{Mit der rechten Seite Aufgenommenes \na Spektrum mit angepassten Gaußkurven}{0.75}
\jafps{rechts/ba_ganz}{Mit der rechten Seite Aufgenommenes \ba Spektrum mit angepassten Gaußkurven}{1}
\jafps{rechts/energy_cal}{}{0.75}
\jaftl{|c|c|c|c|c|}{../data/fits/rechts/na_ganz.table}{naganz}{rechts_na_ganz}
\jaftl{|c|c|c|c|c|}{../data/fits/rechts/ba_ganz.table}{baganz}{rechts_ba_ganz}
\jaft{|c|c|}{rechts_energy_cal.table}{\todo{caption this}}


\todo{gradengleichungen invertieren}
